\chapter{Conclusion}

This project set out to answer a question that every compiler textbook raises and none settles quantitatively: how much does the order of optimization passes matter, and why? The answer was pursued by construction. Rather than probing a production compiler from the outside, the team built a complete compiler infrastructure for a subset of C in which six classical optimizations, constant folding, dead code elimination, common subexpression elimination, copy propagation, strength reduction, and algebraic simplification, are independent, reorderable, and individually correct passes over a three-address-code intermediate representation. Because the pass set is small, the ordering space of 720 permutations could be enumerated in full, and because the compiler is small, every result could be inspected by hand.

The empirical findings are clear. Across fifteen kernels adapted from the MiBench embedded suite and the PolyBench/C numerical suite, the best family of orderings reduces static code size by a geometric mean of 35.7\%, with per-program reductions from 22.3\% to 60.0\%, and reduces the estimated cycle cost by 20.5\% and the dynamic instruction count by 23.2\%. The worst family reduces code size by only 1.3\% and leaves thirteen of the fifteen programs entirely unchanged. Between these extremes lie only nine distinct aggregate outcomes, because the 720 orderings collapse into equivalence classes determined solely by the relative order of four passes: constant folding, copy propagation, common subexpression elimination, and dead code elimination. Strength reduction and algebraic simplification never change the instruction count and affect the cycle estimate only through their position relative to constant folding. The governing rule is simple and was confirmed on every program: the passes that expose opportunities, constant folding and copy propagation, must run before the pass that consumes them, dead code elimination, and dead code elimination should run last. Every one of the 10{,}815 optimized programs produced output identical to its baseline, and every benchmark lowered to RISC-V RV32IM assembly produced the same output in the simulator as in the interpreter, with the code-size improvement measured on the intermediate representation carried through to the assembly.

The engineering findings are equally definite. A compiler infrastructure of about four thousand lines of Python, with a hand-written lexer, a recursive-descent parser, a semantic analyzer, a linear intermediate representation with explicit def-use information, an interpreter, and a stack-slot RISC-V back end, is sufficient to conduct a complete phase-ordering study and is small enough for a student to read in full. The decision to keep every pass local to a basic block and to route all memory access through dedicated opcodes made the passes correct under arbitrary reordering without alias analysis, at a cost in optimization power that the study quantifies. The decision to validate every ordering by execution, and every lowering by simulation, caught three real bugs during development, in common subexpression elimination, in algebraic simplification, and in the interaction heatmap, that static reasoning had missed.

The interactive web application makes the study repeatable by anyone with a browser. A user can load any benchmark or write a new program, inspect every front-end artefact, drag the six passes into any order, watch the intermediate code change stage by stage with a diff and metric deltas, read the RISC-V assembly side by side with the intermediate code that produced it, and run the full exploration to see the Pareto front of size against cycles for their own code. The phase-ordering problem, which is usually described in a paragraph and then set aside, becomes something a learner can manipulate.

The results also situate the project within the literature. The observation that most orderings are redundant is the same observation that made the exhaustive enumerations of Kulkarni and colleagues feasible and that motivates the dormant-pass and dependence-modelling techniques of recent work. The rule that enabling passes must precede clean-up passes is the rule Muchnick stated from experience three decades ago, here given quantitative form on a small pass set. The strong but imperfect correlation between code size and estimated cycles across orderings mirrors the multi-objective framing of POSET-RL and COLE. The project does not claim to improve on any of these; it claims to make their central phenomena exact, visible, and teachable.

\section*{Future Scope}

The infrastructure was designed so that each of the following extensions requires adding to it rather than restructuring it.

\begin{itemize}
  \item \textbf{Global data-flow analysis.} The passes reset their state at every label. Building a control-flow graph and implementing iterative data-flow analysis for reaching definitions, live variables, and available expressions would let the same six passes work across basic blocks, and would raise a new phase-ordering question: whether the structural rules found here survive when the passes become global.
  \item \textbf{More passes and higher-order interactions.} Loop-invariant code motion, induction-variable strength reduction, function inlining, and a peephole pass over the assembly would extend the pass set to ten or more, at which point exhaustive enumeration becomes impractical and the infrastructure can serve as a test bed for the search and learning methods surveyed in Chapter 2, with the exhaustive six-pass results as ground truth.
  \item \textbf{Combined passes.} Following Click and Cooper, merging constant folding, copy propagation, and dead code elimination into a single sparse conditional constant propagation pass would test how much of the best ordering's benefit can be captured by one pass, dissolving part of the ordering problem rather than solving it.
  \item \textbf{Measured rather than estimated cost.} Running the generated assembly under a cycle-accurate simulator, or on a RISC-V development board with a cycle counter, would replace the static cost model with measurements and would allow energy to be studied directly.
  \item \textbf{Register allocation.} A linear-scan or graph-colouring allocator in the back end would remove most of the loads and stores in the generated assembly and would make the assembly-level effect of each ordering more pronounced.
  \item \textbf{Larger and real benchmarks.} Supporting more of C, in particular pointers, structures, and multi-dimensional arrays, would allow unmodified MiBench and PolyBench sources to be compiled, removing the adaptation step.
  \item \textbf{Automated ordering selection.} The equivalence-class structure found here suggests a cheap heuristic: compute the program's instruction-category profile, predict which class it falls in, and choose a representative ordering. Training such a predictor on many programs, and comparing it with the reinforcement-learning approaches in the literature, is a natural next research project.
  \item \textbf{Classroom deployment and evaluation.} The web application should be used in a compiler design course and its effect on students' understanding of pass interactions measured, in the manner of the VisOpt study.
\end{itemize}

With these extensions, the infrastructure can grow from a complete study of six passes into a general platform for phase-ordering research and teaching, while retaining the property that made this study possible: every component is small enough to read, and every result is small enough to check.
